\subsubsection{波动方程的基本解}\label{subsec:wave-fourier-fundamental}

Fourier 变换显然有助于求基本解. 若 $L(D)$ 为常系数算子, 则方程 $L(D)u=\delta$ 变换为 $L(i\bm{\xi})\widehat u=(2\pi)^{-m/2}$, 即一个纯代数方程. 立即得到
\begin{equation}\label{eq:fourier-fundamental-solution-symbol-reciprocal}
\widehat u(\bm{\xi})=\dfrac1{(2\pi)^{m/2}L(i\bm{\xi})},
\end{equation}
唯一的问题是 $L(i\bm{\xi})$ 可能有零点. 若 \eqref{eq:fourier-fundamental-solution-symbol-reciprocal} 有不可积奇性, 则必须考虑适当的正则化. 最后, 还必须计算 $\widehat u(\bm{\xi})$ 的 Fourier 逆变换; 这一步不一定容易.

类似地, Fourier 变换可用于求初值问题的基本解; 现在对 $\mathbb{R}^3$ 中的波动方程这样做. 对问题
\begin{equation}\label{eq:wave-three-dimensional-fundamental-initial-problem}
G_{tt}=\Delta G,\quad G(\bm{x},0)=0,\quad G_t(\bm{x},0)=\delta(\bm{x})
\end{equation}
仅关于空间变量作 Fourier 变换; 即定义
\begin{equation}\label{eq:wave-spatial-fourier-transform}
\widehat G(\bm{\xi},t)=(2\pi)^{-m/2}\int_{\mathbb{R}^m}e^{-i\bm{\xi}\cdot\bm{x}}G(\bm{x},t)\,\mathrm{d}\bm{x},
\end{equation}
并对 \eqref{eq:wave-three-dimensional-fundamental-initial-problem} 作相同类型的变换. 结果为关于变量 $t$ 的常微分方程,
\begin{equation}\label{eq:wave-fourier-transformed-oscillator-problem}
\widehat G_{tt}(\bm{\xi},t)=-|\bm{\xi}|^2\widehat G(\bm{\xi},t),\quad\widehat G(\bm{\xi},0)=0,\quad\widehat G_t(\bm{\xi},0)=(2\pi)^{-3/2}.
\end{equation}
令 $|\bm{\xi}|=\rho$, 容易得到解为
\begin{equation}\label{eq:wave-fourier-oscillator-solution}
\widehat G(\bm{\xi},t)=(2\pi)^{-3/2}\dfrac{\sin\rho t}\rho.
\end{equation}
利用上面的 Example~\ref{ex:fourier-uniform-sphere-mass}, 得到
\begin{equation}\label{eq:wave-three-dimensional-spherical-delta}
G(\bm{x},t)=\dfrac{\delta(r-t)}{4\pi t}.
\end{equation}


\PDEIndexMark{idx:wave-even-dimensional-descent-method}可以证明, 在任意大于 1 的奇数空间维数中, 波动方程的基本解均可用 $\delta(r-t)$ 的导数表示; 由于求解三维以上的波动方程在应用中意义不大, 这里不予证明. 然而, 求解二维波动方程是有意义的. 在偶数空间维数中, 不容易直接计算 $\sin\rho t/\rho$ 的 Fourier 逆变换; 而是使用称为降维法的技巧. 这一技巧基于一个简单观察: 任意二维波动方程的解, 只需与常函数 1 取直积, 就可视为三维波动方程的解. 因此, 可将 \eqref{eq:wave-three-dimensional-spherical-delta} 与 $\delta(x)\delta(y)1(z)$ 作卷积, 得到二维基本解. 利用卷积的定义 \eqref{eq:distribution-convolution-formal-definition}, 计算得到


\begin{equation}\label{eq:wave-descent-convolution-pairing}
\begin{aligned}
&\left(\delta(x)\delta(y)1(z)\dfrac{\delta(r'-t)}{4\pi t},\phi(\bm{x}+\bm{x}')\right)\\
&\quad=\dfrac1{4\pi t}\int_{-\infty}^\infty\int_{r'=t}\phi(x',y',z'+z)\,\mathrm{d}S'\,\mathrm{d}z.
\end{aligned}
\end{equation}
用 $\psi(x,y)$ 表示 $\int_{-\infty}^\infty\phi(x,y,z)\,\mathrm{d}z$, 则 \eqref{eq:wave-descent-convolution-pairing} 简化为
\begin{equation}\label{eq:wave-descent-surface-pairing}
\dfrac1{4\pi t}\int_{r'=t}\psi(x',y')\,\mathrm{d}S',
\end{equation}
计算这一积分得到
\begin{equation}\label{eq:wave-descent-disk-pairing}
\dfrac1{2\pi}\int_{x^2+y^2\leq t^2}\dfrac{\psi(x,y)}{\sqrt{t^2-x^2-y^2}}\,\mathrm{d}x\,\mathrm{d}y.
\end{equation}
由此得到以下二维空间中的基本解:
\begin{equation}\label{eq:wave-two-dimensional-fundamental-solution}
G(\bm{x},t)=\begin{cases}1/2\pi\sqrt{t^2-r^2},&r<t\\0,&r\geq t.\end{cases}
\end{equation}


注意, 热方程基本解的定性性质并不随空间维数真正改变, 但波动方程的基本解却变化显著. 在任意维数中, 波动方程基本解的支集均包含于 $|\bm{x}|\leq t$, 但除此之外, 基本解的形态很不相同. 三维基本解仅支撑于球面 $|\bm{x}|=t$, 而 \eqref{eq:wave-two-dimensional-fundamental-solution} 的支集填满整个圆. Abbott 的 \emph{Flatland}~\cite{Ab} 中的电视机必须设计得与我们的很不相同; 关于这一点, 也见~\cite{Mo}.
